% 29 Mar: clean up

\section{The Position}

John  O. Evjen

\subsection{Installment 1: “Folkebladet” March 13, 1918}


% The Congregation as the Church’s Proper Form

For several years certain disagreements have made themselves felt in Folkebladet concerning the Free Church’s Guiding Principles, regarding what they mean and what consequences they entail. There has been much discussion about it.

The discussion does not seem likely to cease. It is almost as if the Free Church has a broken arm that must be set. The sufferer’s thoughts remain on the broken arm until it is healed. Since the understandings of the principles are exceedingly different, frictions arise. They may often look worse than they are, and those involved in the frictions may sometimes be subjected to harsh judgments. These judgments can sound quite peculiar, for example like the one a little schoolboy wrote in a composition about a horse: “The horse is a very fierce animal. If you attack him he will kick you.” So it can happen that men such as Birkeland and myself may be thought both brute beasts and ‘fierce’.

But, as has been said, so long as the arm is in pain, the thoughts remain fixed on it. And so it is with the principles. It is possible that they mean more than an arm on the Free Church’s body; it is possible that they mean less. At any rate, they play an exceedingly important role in determining the right to vote at the Free Church’s annual meetings. And surely the meaning is that the principles are to be respected between the annual meetings as well. But so long as there prevails lack of clarity with respect to their meaning and their consequences, it is natural that there should be discussion about them.

More recently, the discussion has centered most sharply on the question: Is Folkebladet against Augsburg? In other words, does the present editorial board of Folkebladet interpret the understanding of the congregation maintained in the Guiding Principles, which were in fact essentially drafted by Sverdrup and Oftedal, who were the leading voices both at Augsburg and in Folkebladet?

From the very beginning of his work in America, Folkebladet’s founder,\footnote{Oftedal founded Folkebladet in FIXME.} Prof. Oftedal upheld the concept of the church advanced in the Guiding Principles, while Sverdrup, it seems, did not fully come over to his side until after 1876. They were always agreed about ecclesiastical freedom. But in the first years at Augsburg Sverdrup taught, as we still so often find among Lutheran theologians, that the congregation began at Christ’s ascension, not first at the outpouring of the Spirit. In this connection he could say that it was the “first congregation” that chose an apostle to take Judas’s office, whereas later he always taught that the choice of this apostle was not a congregational election, since the congregation was first born at the outpouring of the Spirit. Likewise, he taught that a local congregation, a church body, a branch of the church, is a member of Christ’s body, a stone in God’s temple. “Christ is the head of all, and every congregation is a member of His body, a living stone in the new temple of God.” A church body is “an independent member of the body.” “The Church as a whole is Christ’s body.” (Collected Writings I.)\footnote{Professor Georg Sverdrup Biblical and Church-Historical Sketches and Addresses Edited by Andreas Helland 1909)}


% Potential footnote V1 page 83:

%Men for at et Samfund skal have Ret til at bestaa ved Siden af andre, skal have Ret til at være et selvstændigt Lem paa Legemet, da maa det have sin særegne Opgave, sin eiendommelige Virksomhed. Det er derom der vil blive Spørgsmaal, og ethvert Samfund, som ikke er istand til at eftervise sin særegne Opgave i den kirkelige Udvikling, det mangler Grundbetingelsen for at være selvstændigt. Men hvor et Kirkesamfund kan eftervise, at det har en Gjerning at gjøre, som vilde forsømmes, om det ikke var, at det har en Side af Kristi Er kjendelse at fremholde og en Gren af Kristi Liv at udfolde, som ikke finder Plads indenfor de allerede bestaaende Kirkesamfund, der hardet sin Ret til at bestaa, og der har de allerede bestaaende Samfund

After 1876, and right up to his death, Sverdrup taught, like Oftedal, that the local congregation—that is, the believers within it—is not a stone in the temple, nor a member of Christ’s body, but is itself the temple, God’s house, Christ’s body. The other ecclesiastical organizations he regarded as historical arrangements, not divinely instituted in the way the local congregation is, though useful organs for the congregation. Both Sverdrup and Oftedal were, until the 1890s, advocates of church association in the sense that they desired an association of congregations; but they were not church-body men in the sense that the church body or synod was anything more than an outward expedient.

% Potential footnote V1 Page 95
% Christ is the Head of all, and every congregation a member of His body, a living stone in the new temple of God

% mange Menigheder er i en Forstand ikke mange, men kun een. Thi naar Paulus (Ap. Gj. 20, 28) formaner de Ældste i Efesus, at de skal vogte Guds Menighed, som han har forhvervet sig med sit Blod, da er det vel neppe kun Menigheden i Efesus, som saaledes er Guds Eiendomsfolk, men det er den hele kristne Menighed, hvor den end findes. Og det er i Virkeligheden saa, at dersom Kristus er levende tilstede i Ord og Sakrament, og hver Menighed kun derfor er en kristen Menighed, fordi den er bleven delagtig i dette Kristi Liv, saa er jo Kristus alles Hoved og enhver Menighed et Lem paa hans Legeme, en levende Sten i det nye Guds Tempel. Ordet, Daaben og Nad veren er ligesom de tre Baand, der knytter de enkelte Menigheder sammen, netop fordi de er de tre Midler, hvorved Kristi Liv kan med deles til den faldne Verden, og derfor forplante den nye Menneske slegt, hvis levendegjørende Aand er Kristus. Det nye Testamente af lægger Vidnesbyrd om, at disse tre er det levende Enhedsbaand mel lem alle kristne og alle Menigheder, og det kjender ikke noget andet. Det er om Ordets Virkning Apostelen Paulus taler, naar han

In faculty opinions, treatises, and in part also sermons, the theological professors at Augsburg again and again set forth that the congregation (the local congregation) is the “only outward form of the Church” (the Church—the holy catholic Church; compare question 243 in the Free Church’s edition of Pontoppidan’s Explanation): “the proper divine order within the Christian religion”; “the truly proper form of Christianity’s manifestation in the world”; “the proper form of God’s kingdom on earth”; the only outward form of His kingdom on earth appointed by God; “the foundational truth of the Church is and remains according to God’s Word the congregation alone and no church body”; “according to Scripture the congregation is the highest form of association”; “the outward form of the kingdom of God, which is the congregation.” Prof. Oftedal wrote in the 1870s that Elling Eielsen abolished the visible Church in order to set the invisible Church, the communion of saints, in its place. In Eielsen’s understanding he found “no room for the only outward form of the Church that God’s Word mentions, namely the congregation.” In a long joint declaration from the 1870s it was stated that the church body “has no such divine right as the congregation, but, like all outward church bodies, has arisen through historical development”; likewise that “it belongs inseparably to the essence of the Christian Church to exist in the form of congregations.”

That it is local congregations everywhere that are meant here is self-evident. But this will, moreover, also appear from many other citations in this article. The congregation is everywhere presented as the proper form, set by God, for the holy catholic Church, the Church, or the kingdom, in contrast to the state-church form, the church-body form, the papal-church form, and the ‘free-free’ form. For the time being, it will suffice to refer to the treatise “What Is the Congregation” in Collected Writings II, p. 12, where we read this:

\begin{quote}
Our children’s catechism teaches what the Church is, but not directly of what the congregation is. The New Testament, however, speaks all the more of the congregation or the congregations. Now people sometimes call them local congregations, or ‘local churches.’ It is by no means free from a certain tinge of disparagement when people speak of ‘local congregations.’ We therefore take the liberty of observing that what is spoken of here is the ‘local congregation.’ And since it is an indisputable fact that the New Testament speaks with the deepest reverence and utmost holy seriousness precisely of the local congregations, and calls them God’s congregations and God’s temple and Christ’s body and Christ’s bride, it is not too much if we also, with all reverence and respect, speak of this divine institution among us and esteem it very highly for Christ’s and God’s sake.
\end{quote}


Let it be noted that the local congregation is designated a ‘divine institution.’

Sverdrup considered it a deficiency in the children’s catechism that it did not speak directly of what the congregation was. For this reason he altered question 243\footnote{FIXME to point to question 243} in the Explanation, and wrote: “Is the holy catholic Church found here among us?” His answer reads: “Yes, in the congregation, which in Scripture is called Jesus Christ’s body and God’s house.” The next question, concerning the outward congregation, might lead one to think that the word “congregation” in the answer to question 243 means all believers on earth. But that would be a hasty conclusion, which cannot be defended in the light of all that both Sverdrup and Oftedal wrote about “congregation” and “outward congregation.”

Thus the congregation is the proper form. It is governed by the law of Christ’s body. 

Church bodies, state churches, and free churches (associations of congregations) are not governed by the law of Christ’s body (Guidance II, 33, 145).\footnote{FIXME point to Guidance II, 33, 145} We may use an analogy: the congregation is the proper form of the kingdom of God, somewhat as monogamy is the highest form of union between man and woman, in contrast to polygamy. This is admittedly a very imperfect comparison; it limps, as all figures do, especially if one were to regard the other forms of the kingdom as somehow immoral. But the point is plain enough.

This understanding of the congregation as the “proper form” and Christ’s body is, so far as Lutherans are concerned, distinctive of the Free Church (and we may add the free-church people in Norway, for example Odland, Jens Greve, and Solem; but not The Evangelical Lutheran Free Church). Otherwise I know of no other Lutheran church body in the world that maintains this understanding. All other Lutheran church bodies maintain that “Christ’s body,” when applied to the concept of the Church, has reference to the holy catholic Church, which consists of all believers on earth, but has no reference to any local congregation. In other words, “Christ’s body” must be applied to the whole organism, not to any organization.


In this respect, the Free Church stands almost alone among Lutherans. It is true that a couple of editors in the new church body, The Norwegian Lutheran Church body, have stated that this church body also teaches that the congregation is the “proper form.” But their statements show that they have not penetrated deeply enough into the Free Church understanding of the congregation. And one cannot hold that against them, seeing that so many of our own people do not have one and the same understanding of principle no. 1 in our “Guiding Principles.” Besides, they are contradicted by Pastor Langjærd in “The Lutheranism of Our Four-Hundred-Year Lutheranism.”

But if this principle is peculiar to the Free Church as a Lutheran ‘body,’ then it shares that principle with Congregationalists and Baptists. It may indeed offend many to hear that we are in such company. But a fact is not abolished by offense. Nor does it make Sverdrup and Oftedal dependent on the Reformed, any more than The Open Declaration—simply because it accused the Synod of Catholicism—made Synod people into Catholics, so it does not make Sverdrup and Oftedal dependent upon the Reformed, although Sverdrup already in 1876 thought that we could learn from “such a congregational church as the Reformed has long been,” and that “it would be sheer self-conceit to think that there was nothing to learn from those who have tried these things before us” (the Reformed in America). Collected Writings I, 143.

One cannot expect the ordinary man, nor even the pastor strongly occupied in practical work, to be able to find time and occasion to acquaint himself very much with the Lutheran and Reformed literature on church polity. But it is fair to expect of one who writes publicly about the Guiding Principles and criticizes others’ understanding of them, that he acquaint himself with them well enough to know what the principles mean, what is consistent with them, and how much they have in common with Luther and how much or how little they have in common with Calvin. One does not gain much insight into the ‘Guiding Principles’ by reading dogmatics, for dogmatics treats the congregation as an object of faith. But one finds light in the literature on church polity and in church history.

But do we not already have all the teaching we need? Have we not sat at the feet of the “older teachers”?

Yes—but how much did we hear there about the “Guiding Principles” and the Augsburg understanding of the congregation? From Oftedal, almost nothing; from Sverdrup, a few lectures, given every third year when he came to the end of Lønning’s dogmatics, which he mercilessly tore to pieces.

But what we did not hear, we did get to read in the papers. Several of us have read Sverdrup’s and Oftedal’s articles again and again, and in context. Here at Augsburg almost everything that Sverdrup and Oftedal wrote is accessible; and almost all the books they used as professors (including the marginal notes they made) are accessible. As far as possible I have tried to acquaint myself with their work as teachers and public writers, so that I dare say that I have a fairly well-grounded knowledge of their ecclesiastical views.

I hope that no one will take this as though it were spoken in a proud spirit. I am well aware of my limitations. As for me, I did not gallop through Augsburg, as many others did. And when I was finished with Augsburg I had to continue my schooling for three to four years before I could, with a reasonably good scholarly conscience—if one may use that expression, present myself as a candidate for a pastoral call or a teaching post. And finally, after several years of teaching at other schools, I had thus come to know Augsburg both at close hand and from afar, and now once again I came to know it more intimately than ever before. (To be continued.)

